\section{Score Matching：学习 Score Function}

\subsection{问题设定}

前面我们看到：（1）score function 完全描述了概率分布；（2）通过 Langevin dynamics 可以利用 score function 采样。那么关键问题变成了：\textbf{如何从数据中学习 score function？}

最直接的想法是定义一个 score 匹配目标：

$$\mathcal{L}_{\text{SM}} = \mathbb{E}_{q(x)}\left[\left\| s_\theta(x) - \nabla_x \log q(x) \right\|^2\right]$$

其中 $s_\theta(x)$ 是我们用神经网络参数化的 score estimator。

\begin{warningbox}{直接 Score Matching 的困难}
上面的损失函数有一个根本性的问题：它需要知道真实分布 $q(x)$ 的 score function $\nabla_x \log q(x)$，而这正是我们要学习的目标。这似乎陷入了``鸡生蛋、蛋生鸡''的循环。
\end{warningbox}

\subsection{Denoising Score Matching}

解决这个问题的关键洞察是：虽然我们不知道原始数据分布 $q(x)$ 的 score function，但如果我们定义一个\textbf{加噪版本}的分布 $q_\sigma(x_t)$：

$$q_\sigma(x_t) = \int q(x_t \mid x_0)\, q(x_0)\, dx_0$$

其中 $q(x_t \mid x_0) = \mathcal{N}(x_t;\; x_0, \sigma^2 I)$，那么\textbf{条件 score} $\nabla_{x_t} \log q(x_t \mid x_0)$ 是已知的：

$$\nabla_{x_t} \log q(x_t \mid x_0) = -\frac{x_t - x_0}{\sigma^2}$$

\begin{importantbox}{Denoising Score Matching (DSM) 损失}
经过数学推导（省略恒等变换的细节），可以证明 score matching 目标等价于：
$$\mathcal{L}_{\text{DSM}} = \mathbb{E}_{q(x_0)}\, \mathbb{E}_{q(x_t \mid x_0)}\left[\left\| s_\theta(x_t) - \nabla_{x_t} \log q(x_t \mid x_0) \right\|^2\right]$$
由于 $\nabla_{x_t} \log q(x_t \mid x_0)$ 有闭式解，这个损失是可计算的！
\end{importantbox}

\begin{knowledgebox}{DSM 的训练过程}
Denoising Score Matching 的训练步骤极其简单：
\begin{enumerate}
    \item 采样一个数据点 $x_0 \sim q(x_0)$（即取一个训练样本）。
    \item 采样噪声 $\epsilon \sim \mathcal{N}(0, I)$，得到 $x_t = x_0 + \sigma\, \epsilon$。
    \item 计算 score 目标 $\nabla_{x_t} \log q(x_t \mid x_0) = -\epsilon / \sigma$。
    \item 最小化 $\|s_\theta(x_t) + \epsilon/\sigma\|^2$。
\end{enumerate}
可以看到，这与 DDPM 的噪声预测训练在本质上是\textbf{完全相同}的。
\end{knowledgebox}

\subsection{与 DDPM 训练的等价性}

在 DDPM 中，噪声预测器的训练目标是：

$$\mathcal{L}_{\text{DDPM}} = \mathbb{E}_{x_0, t, \epsilon}\left[\left\| \epsilon_\theta(x_t, t) - \epsilon \right\|^2\right]$$

而在 denoising score matching 中，score 网络的训练目标是：

$$\mathcal{L}_{\text{DSM}} = \mathbb{E}_{x_0, t, \epsilon}\left[\left\| s_\theta(x_t, t) + \frac{\epsilon}{\sqrt{1-\bar{\alpha}_t}} \right\|^2\right]$$

两者的关系是 $s_\theta(x_t, t) = -\epsilon_\theta(x_t, t) / \sqrt{1-\bar{\alpha}_t}$，即 score 网络和噪声预测网络只差一个与时间步相关的缩放因子。

\begin{importantbox}{DDPM 与 Score-Based Models 的统一}
DDPM 的噪声预测框架和 score-based 的 score matching 框架是\textbf{同一枚硬币的两面}：
\begin{itemize}
    \item DDPM 视角：训练一个神经网络来预测加在数据上的噪声。
    \item Score-based 视角：训练一个神经网络来估计加噪数据分布的 score function。
    \item 两种视角在数学上完全等价，只差一个缩放因子。
\end{itemize}
\end{importantbox}

\subsection{本章小结}

Score matching 通过巧妙地将不可计算的边缘 score matching 目标转化为可计算的 denoising score matching 目标，解决了``如何从数据中学习 score function''的问题。这个训练目标与 DDPM 的噪声预测训练在数学上完全等价，从而将两个看似不同的研究方向统一起来。

